\section{The sharp global single-parity bound}
\label{app:global-single-parity}

Write
\[
 C_n(u)=\E\chi_{[n]}(X)\tanh(u\cdot X),\qquad
 L(a)=\sqrt{2/\pi}\,\frac{F_0(a)}a,
 \qquad L_* = L(a_0).
\]
Here $a_0\in(1,2)$ is the limiting balanced gain from
Theorem~\ref{thm:curvature}. It is the unique global maximizer of $L$.
Here the strictness survives the limit quantitatively. For $a\le1$,
the positive-series proof gives
\[
 a(q_n/a)'\ge2\delta\,q_n/a,
 \qquad \delta=44428/44251-1>0.
\]
Divide by $\kappa_n$ and pass to the locally uniform limit to obtain
$a(F_0/a)'\ge2\delta F_0/a>0$.
For $a\ge1$, the alternating-series proof of concavity has a
uniformly positive first-pair gap: with
$p(t)=-\tanh''t>0$, its normalized limit obeys
\[
 F_0''(a)\le-2[p(a)-9p(3a)]<0.
\]
The inequality $9p(3a)<p(a)$ is the explicit ratio bound in
Theorem~\ref{thm:local}. Thus $aF_0'-F_0$ is strictly decreasing
above one. Its positive sign at one and negative sign at two also
follow from the strict normalized endpoint bounds in that theorem.
There is exactly one zero $a_0\in(1,2)$, and the limiting scalar
maximum is nondegenerate. We first locate every
asymptotically optimal unrestricted neuron.

\begin{lemma}[Asymptotic localization at one gain]
\label{lem:asymptotic-single-parity}
Along odd integers $n\to\infty$,
\[
 n\sup_{u\ne0}\frac{|C_n(u)|}{\|u\|_2}\longrightarrow L_*.
\]
Every sequence with scaled gain at least $L_*-o(1)$ satisfies, after
input sign changes,
\[
 \frac1n\#\{i:|u_i-a_0|>\epsilon\}\longrightarrow0
 \quad(\epsilon>0),\qquad
 \frac{\|u\|_2^2}{n}\longrightarrow a_0^2.
\]
After taking a subsequence there are $A_n\to a_0$ and sets $I_n$ of
size $n-o(n)$ such that
\begin{equation}\label{eq:single-parity-bulk}
 \max_{i\in I_n}|u_i-A_n|\to0,\qquad
 \sum_{i\in I_n}(u_i-A_n)^2\le K,\qquad
 \sum_{i\notin I_n}u_i^2=o(n).
\end{equation}
The constant $K$ is absolute for any fixed positive competitive-gain
threshold below $L_*$.
\end{lemma}
\begin{proof}
Balanced neurons give the lower bound. Consider any sequence with
scaled gain bounded below by a positive constant. By
Theorem~\ref{thm:odd-lattice-localization}, after input sign changes,
\[
 u_i=a_n r_i+d_i,\qquad r_i\in\{1,3,5,\ldots\},\qquad
 a_n\in[a_-,a_+],\qquad
 \sum_i d_i^2\le K,\quad \sum_i r_i^2\le K_1 n.
\]
Pass to a subsequence with $a_n\to a$, with empirical probabilities
of each odd integer converging to $p_r$, and with
$n^{-1}\sum_i r_i^2\to Q_{\rm total}$. The second-moment bound makes
the empirical laws tight, so $\sum_rp_r=1$. Put
\[
 Q_{\rm lim}=\sum_rp_r r^2\le Q_{\rm total}.
\]
The bounded residual energy implies
$\|u\|_2^2/n\to a^2Q_{\rm total}$.

Fix a finite cutoff $J$ containing a type of positive limiting
frequency. Retain the coordinates of types $r\le J$ with $p_r>0$,
and discard $o(n)$ further coordinates so their maximum residual
tends to zero. This is possible by deleting $|d_i|>\epsilon_n$,
where $\epsilon_n\downarrow0$ sufficiently slowly that
$K/\epsilon_n^2=o(n)$. Let $g_J$ be the gcd of the retained types and
$Q_J=\sum_{r\le J}p_r r^2$. Least-squares refitting their common scale
changes $a_ng_J$ by $O(n^{-1/2})$, decreases residual energy, and
preserves the vanishing maximum residual. Condition on every other
selected coordinate. The uniform-in-bias residual central limit
theorem, Theorem~\ref{thm:residual-clt}, and
$\|T_{A,E}\|_\infty\le F_0(A)$ give
\[
 |C_n(u)|\le
 \sqrt{\frac2\pi}\,
 \frac{F_0(ag_J)}{\sqrt{nQ_J/g_J^2}}+o(n^{-1/2}).
\]
The outside parity sign does not affect this upper bound. Therefore
\begin{equation}\label{eq:lattice-truncation-gain}
 \limsup_n\frac{n|C_n(u)|}{\|u\|_2}
 \le L(ag_J)\frac{g_J^2}{\sqrt{Q_JQ_{\rm total}}}.
\end{equation}
As $J$ increases through the support of $p$, its positive integer
gcds eventually stabilize at a value $g$. Also $Q_J\uparrow Q_{\rm lim}$.
Every support point is a positive odd multiple of $g$, whence
$Q_{\rm total}\ge Q_{\rm lim}\ge g^2$. Equation
\eqref{eq:lattice-truncation-gain} proves the upper bound $L_*$.
Equality at this bound forces
\[
 ag=a_0,\qquad Q_{\rm total}=Q_{\rm lim}=g^2.
\]
Thus $p_g=1$ and no second-moment mass escapes to larger harmonics.
Retain the coordinates of type $g$ with $|d_i|\le\epsilon_n$.
They give \eqref{eq:single-parity-bulk} with $A_n=a_ng$.
Applying this argument to every subsequence proves the asserted
empirical and norm conclusions for the original sequence.
\end{proof}

\begin{lemma}[An exactly equal majority at a global maximizer]
\label{lem:exact-majority}
For any sequence of global maximizers of $|C_n(u)|/\|u\|_2$ with
odd $n\to\infty$, after input sign changes there is an exactly equal
bulk of size $n-o(n)$, of gain tending to $a_0$. All remaining
coordinates have total squared norm $o(n)$. Moreover, one fixed
neighborhood of the common bulk gain contains no exceptional
coordinate for all sufficiently large $n$.
\end{lemma}
\begin{proof}
For each fixed $n\ge3$ the ratio tends to zero at infinity. It also
tends to zero at the origin, since the top Fourier coefficient
vanishes to order at least $n$ there. Thus a nonzero global maximum
is attained. None of its selected coordinates is zero. Apply
Lemma~\ref{lem:asymptotic-single-parity} and use its set $I_n$.

Choose $i,j\in I_n$, and put
$\mu=(u_i+u_j)/2$ and $h=(u_i-u_j)/2$. The remaining selected
degree $n-2$ is odd, so its response
\[
 G(b)=\E\chi_{[n]\setminus\{i,j\}}(X)
       \tanh\left(b+\sum_{l\ne i,j}u_lX_l\right)
\]
is even. Conditioning the two removed signs gives the exact identity
\begin{equation}\label{eq:pair-response}
 C_n(u)=\tfrac12\{G(2\mu)-G(2h)\}.
\end{equation}
Remove this pair from $I_n$ and least-squares refit its common gain
$a_n$. Its residual energy stays bounded, its maximum residual tends
to zero, and $a_n\to a_0$. Let $N=n-o(n)$ be the remaining bulk size,
let $E$ be its residual energy, and write $W$ for the sum of the
other selected coordinates with their independent signs. Then
$\E W^2=o(n)$. Theorem~\ref{thm:residual-clt}, including two bias
derivatives, gives
\[
 G(b)=\epsilon_N\sqrt{\frac2{\pi N}}
 \E\left[\chi_{\rm out}
 e^{-(b+W)^2/(2a_n^2N)}
 T_{a_n,E}(b+W+\eta a_n)\right]+o(n^{-1/2}),
 \qquad \eta=N\bmod2.
\]
The statement and its first two derivatives are uniform in $b$ on
compact sets and uniform over the removed pair. The Gaussian
envelope can be removed in expectation there, since
$\E W^2=o(n)$; its first and second derivative contributions also
tend to zero after multiplication by $\sqrt n$. The periodic
derivatives are uniformly bounded. Thus, for $j=0,1,2$,
\begin{equation}\label{eq:pair-clt-derivatives}
 G^{(j)}(b)=\epsilon_N\sqrt{\frac2{\pi N}}
 \E[\chi_{\rm out}T_{a_n,E}^{(j)}(b+W+\eta a_n)]
 +o(n^{-1/2}).
\end{equation}
Every parity sign can be absorbed into a shift by $2a_n$, using
anti-periodicity. Put $s=\operatorname{sign}C_n(u)$ and $H=-sG$.
Consequently \eqref{eq:pair-clt-derivatives} represents $H$ and its
derivatives as probability averages of the same smoothed periodic
response, with positive normalizing factor $\sqrt{2/(\pi N)}$.

Since $\mu-a_n\to0$ and $h\to0$,
anti-periodicity gives $G(2\mu)=-G(0)+o(n^{-1/2})$ and
$G(2h)=G(0)+o(n^{-1/2})$. Global optimality and
\eqref{eq:pair-response} therefore force the normalized average
representing $H(0)$ to tend to $F_0(a_0)$. Choose a fixed small
$\delta>0$. Uniform boundedness of the first periodic derivative
then makes this average at least $4F_0(a_n)/5$ throughout
$|b|\le\delta$, for all large $n$ and uniformly in the pair.

Lemma~\ref{lem:periodic-curvature} is preserved by Gaussian smoothing
and probability averaging. It and \eqref{eq:pair-clt-derivatives}
therefore give one $c>0$ with
\[
 H''(b)\le-c/\sqrt n\qquad(|b|\le\delta).
\]
Keep the pair mean fixed and vary its half-difference $x$ toward
zero. Its positive signed coefficient is
$q(x)=\tfrac12\{H(2x)-H(2\mu)\}$. Evenness gives $q'(0)=0$,
and $q''(x)=2H''(2x)<0$ throughout the relevant segment. Unless
$h=0$, balancing the pair strictly increases the coefficient while
strictly decreasing its squared norm
$\|u_{\rm rest}\|_2^2+2\mu^2+2x^2$. This contradicts global
maximality. All coordinates of $I_n$ are therefore exactly equal.

The same proof includes one additional coordinate within a fixed
sufficiently small distance $\delta_0$ of the common gain. Its
pair mean and half-difference change the normalized response by
$O(\delta_0)$; removing it from the outside sum preserves
$\E W^2=o(n)$. Choose $\delta_0$ so that the $4/5$ peak threshold
and the curvature interval still apply. Such a coordinate must
equal the bulk gain too. This proves the final assertion.
\end{proof}

\begin{theorem}[Sharp global affine tanh parity bound]
\label{thm:global-single-parity}
For every sufficiently large odd $n$, and every affine tanh atom on
any finite Boolean cube containing a selected set $A$ of size $n$,
\[
 w_n\left|\E\chi_A(X)\tanh(b+u\cdot X)\right|
 \le\sqrt{b^2+\|u\|_2^2}.
\]
At nonzero equality, the bias and all coordinates outside $A$ vanish,
and the selected coordinate magnitudes all equal the unique optimal
balanced gain $a_n$.
\end{theorem}
\begin{proof}
First consider the bias-free intrinsic problem. Suppose nonbalanced
global maximizers exist for arbitrarily large odd $n$. By
Lemma~\ref{lem:exact-majority}, write their exactly equal bulk as
$N=n-k$ coordinates of gain $a\to a_0$, with $k=o(n)$ and outside
weights $w_i$ satisfying $\sum_iw_i^2=o(n)$. Every exceptional
coordinate is outside one fixed neighborhood of $a$.

Conditioning on the $k$ exceptional signs and applying the equal-bulk
case of Theorem~\ref{thm:residual-clt} gives, with
$B=\sum_iw_iY_i$ and $\eta=N\bmod2$,
\[
 C_n=\epsilon_N\sqrt{\frac2{\pi N}}
 \E\left[\chi_{[k]}(Y)e^{-B^2/(2a^2N)}T_a(B+\eta a)\right]
 +o(n^{-1/2}).
\]
Since $\E B^2=o(n)$ the Gaussian envelope can be removed. Global
competitiveness and norm asymptotics imply
\[
 \left|\E\chi_{[k]}(Y)T_a\left(\eta a+\sum_iw_iY_i\right)\right|
 \longrightarrow F_0(a).
\]
The integers $\eta$ and $k$ have odd sum. The dimension-free phase
stability theorem, Theorem~\ref{thm:phase-stability}, yields
\[
 \sum_i\operatorname{dist}(w_i,(2\mathbb Z+1)a)^2\longrightarrow0.
\]
The fixed-neighborhood conclusion of Lemma~\ref{lem:exact-majority}
excludes the first positive odd multiple for every outlier. Hence
$w_i\ge3a-o(1)$ uniformly over them, and
\begin{equation}\label{eq:exception-norm-inflation}
 \|u\|_2^2\ge a^2\{n+(8-o(1))k\}.
\end{equation}

Put $\delta=0$ if $N$ is odd and $\delta=1$ otherwise.
Theorem~\ref{thm:equal-bulk-bias} gives
\[
 |C_n|\le q_{N+\delta}(a)(1+O(N^{-2})),
\]
where the factor is absent for odd $N$. The balanced scalar expansion,
uniform on $a\in[1,2]$, gives for $h=o(n)$ with $n-h$ odd
\[
 \frac{q_{n-h}(a)}{q_n(a)}
 =\sqrt{\frac n{n-h}}\{1+O(h/n^2)\}.
\]
For $h=0$ the ratio is exactly one. Taking $h=k-\delta$ and using
$k\ge1$, the even-bulk error is absorbed into $O(k/n^2)$. Thus
\[
 \frac{|C_n|/\|u\|_2}{q_n(a)/(a\sqrt n)}
 \le\frac{n}{\sqrt{(n-k+\delta)(n+(8-o(1))k)}}
       \{1+O(k/n^2)\}<1
\]
for all sufficiently large $n$. This contradicts global maximality,
since the denominator is the gain of a balanced competitor. Therefore
every global intrinsic maximizer is balanced, and
Theorem~\ref{thm:local} identifies its gain as $a_n$.

For a fixed bias $B$ and any selected coordinate $i$, the odd-degree
identity \eqref{eq:bias-identity} reads
\[
 C_n(B;u)=\tfrac12\{C_n(0;u+Be_i)+C_n(0;u-Be_i)\}.
\]
The intrinsic bound and square-root concavity imply
$w_n|C_n(B;u)|\le\sqrt{\|u\|_2^2+B^2}$.
Conditioning outside coordinates and applying Jensen proves the stated
affine bound. At a nonzero equality every selected coordinate is
nonzero. The square-root concavity step is strict whenever $B\ne0$,
because the two squared norms differ by $4u_iB$. Conditioning therefore
forces $b+u_{A^c}\cdot X=0$ at every outside input, hence zero bias and
outside weights. Intrinsic equality gives the asserted selected gains.
\end{proof}

The support estimate needed for the code theorem is
Lemma~\ref{lem:absorb-support}. It follows directly from this global
intrinsic inequality and the quantitative square-root concavity gap;
no quantitative stability of the selected balanced weights is needed.
